The proposed PaperSwarm framework demonstrates a significant improvement over the baseline ordinary least squares (OLS) regression method. Specifically, the relative reduction in mean test mean squared error (MSE) achieved by ridge regression over OLS is reported as \result{cl-improve} percent. This improvement underscores the effectiveness of regularization in mitigating overfitting, especially in scenarios with high feature dimensionality and limited training data.

The selected ridge penalty \result{cl-alpha} was found to yield the lowest mean test MSE, highlighting the importance of carefully tuning regularization parameters. The standard deviation of the OLS test MSE across seeds \result{cl-ols-std} indicates high variance, suggesting that OLS solutions can be unstable with different initial conditions. In contrast, the ridge regression model exhibits lower variance, as evidenced by the smaller standard deviation of the test MSE \result{cl-ridge-std}.

The mean test MSE of the minimum-norm OLS solution, averaged over all seeds, is \result{cl-ols-mean}, while the mean test MSE of ridge regression is \result{cl-ridge-mean}. These results confirm that while OLS can achieve a lower mean test MSE in some cases, ridge regression generally provides more stable and reliable performance. The irreducible noise floor, used as a reference for MSE, is \result{cl-noise}, indicating the inherent variability in the data that cannot be reduced by any model.

The experiments were conducted using \result{cl-seeds} random seeds to ensure robustness, with \result{cl-features} features and \result{cl-train} training samples per split. The test set comprised \result{cl-test} samples, allowing for a fair evaluation of model generalization. These findings support the claim that PaperSwarm can effectively generate reproducible and high-quality research papers by leveraging the strengths of multi-agent systems and evidence-gating mechanisms.
